\documentclass[10pt,a4paper]{amsart}
\usepackage{amsmath,amssymb,amsthm,graphicx}
\usepackage{verbatim}
\newenvironment{reference}{\comment}{\endcomment}
%\newenvironment{reference}{}{}
\newenvironment{slogan}{\comment}{\endcomment}
\newenvironment{history}{\comment}{\endcomment}

% For commutative diagrams we use Xy-pic
\usepackage[all]{xy}

% We use 2cell for 2-commutative diagrams.
\xyoption{2cell}
% Theorem environments.
%
\theoremstyle{plain}
\newtheorem{theorem}[subsection]{Theorem}
\newtheorem{proposition}[subsection]{Proposition}
\newtheorem{lemma}[subsection]{Lemma}

\theoremstyle{definition}
\newtheorem{definition}[subsection]{Definition}
\newtheorem{example}[subsection]{Example}
\newtheorem{exercise}[subsection]{Exercise}
\newtheorem{situation}[subsection]{Situation}

\theoremstyle{remark}
\newtheorem{remark}[subsection]{Remark}
\newtheorem{remarks}[subsection]{Remarks}

\numberwithin{equation}{subsection}

% Macros
%
\def\lim{\mathop{\mathrm{lim}}\nolimits}
\def\colim{\mathop{\mathrm{colim}}\nolimits}
\def\Spec{\mathop{\mathrm{Spec}}}
\def\Hom{\mathop{\mathrm{Hom}}\nolimits}
\def\Ext{\mathop{\mathrm{Ext}}\nolimits}
\def\SheafHom{\mathop{\mathcal{H}\!\mathit{om}}\nolimits}
\def\SheafExt{\mathop{\mathcal{E}\!\mathit{xt}}\nolimits}
\def\Sch{\mathit{Sch}}
\def\Mor{\mathop{\mathrm{Mor}}\nolimits}
\def\Ob{\mathop{\mathrm{Ob}}\nolimits}
\def\Sh{\mathop{\mathit{Sh}}\nolimits}
\def\NL{\mathop{N\!L}\nolimits}
\def\CH{\mathop{\mathrm{CH}}\nolimits}
\def\proetale{{pro\text{-}\acute{e}tale}}
\def\etale{{\acute{e}tale}}
\def\QCoh{\mathit{QCoh}}
\def\Ker{\mathop{\mathrm{Ker}}}
\def\Im{\mathop{\mathrm{Im}}}
\def\Coker{\mathop{\mathrm{Coker}}}
\def\Coim{\mathop{\mathrm{Coim}}}

% Boxtimes
%
\DeclareMathSymbol{\boxtimes}{\mathbin}{AMSa}{"02}

%
% Macros for moduli stacks/spaces
%
\def\QCohstack{\mathcal{QC}\!\mathit{oh}}
\def\Cohstack{\mathcal{C}\!\mathit{oh}}
\def\Spacesstack{\mathcal{S}\!\mathit{paces}}
\def\Quotfunctor{\mathrm{Quot}}
\def\Hilbfunctor{\mathrm{Hilb}}
\def\Curvesstack{\mathcal{C}\!\mathit{urves}}
\def\Polarizedstack{\mathcal{P}\!\mathit{olarized}}
\def\Complexesstack{\mathcal{C}\!\mathit{omplexes}}
% \Pic is the operator that assigns to X its picard group, usage \Pic(X)
% \Picardstack_{X/B} denotes the Picard stack of X over B
% \Picardfunctor_{X/B} denotes the Picard functor of X over B
\def\Pic{\mathop{\mathrm{Pic}}\nolimits}
\def\Picardstack{\mathcal{P}\!\mathit{ic}}
\def\Picardfunctor{\mathrm{Pic}}
\def\Deformationcategory{\mathcal{D}\!\mathit{ef}}
\usepackage{cleveref}
\crefname{originalequation}{original equation}{original equations}
\Crefname{originalequation}{Original equation}{Original equations}
\makeatletter
\newlabel{equation-base-change}{{0A9J}{source}{}{sourcepointer-1}{}}
\newlabel{equation-base-change@cref}{{[originalequation][0][]0A9J}{[1][0][]source}{}{}{}}
\makeatother
\crefname{originalequation}{original equation}{original equations}
\Crefname{originalequation}{Original equation}{Original equations}
\makeatletter
\newlabel{equation-base-change-map}{{0AA6}{source}{}{sourcepointer-2}{}}
\newlabel{equation-base-change-map@cref}{{[originalequation][0][]0AA6}{[1][0][]source}{}{}{}}
\makeatother
\crefname{originaltheorem}{original theorem}{original theorems}
\Crefname{originaltheorem}{Original theorem}{Original theorems}
\makeatletter
\newlabel{perfect-theorem-bondal-van-den-Bergh}{{09IS}{source}{}{sourcepointer-3}{}}
\newlabel{perfect-theorem-bondal-van-den-Bergh@cref}{{[originaltheorem][0][]09IS}{[1][0][]source}{}{}{}}
\makeatother
\crefname{originalreference}{original reference}{original references}
\Crefname{originalreference}{Original reference}{Original references}
\makeatletter
\newlabel{perfect-remark-pullback-generator}{{0BQT}{source}{}{sourcepointer-4}{}}
\newlabel{perfect-remark-pullback-generator@cref}{{[originalreference][0][]0BQT}{[1][0][]source}{}{}{}}
\makeatother
\crefname{originalsection}{original section}{original sections}
\Crefname{originalsection}{Original section}{Original sections}
\makeatletter
\newlabel{cohomology-section-global-RHom}{{0B6A}{source}{}{sourcepointer-5}{}}
\newlabel{cohomology-section-global-RHom@cref}{{[originalsection][0][]0B6A}{[1][0][]source}{}{}{}}
\makeatother
\crefname{originallemma}{original lemma}{original lemmas}
\Crefname{originallemma}{Original lemma}{Original lemmas}
\makeatletter
\newlabel{perfect-lemma-affine-morphism-and-hom-out-of-perfect}{{0AA7}{source}{}{sourcepointer-6}{}}
\newlabel{perfect-lemma-affine-morphism-and-hom-out-of-perfect@cref}{{[originallemma][0][]0AA7}{[1][0][]source}{}{}{}}
\makeatother
\crefname{originallemma}{original lemma}{original lemmas}
\Crefname{originallemma}{Original lemma}{Original lemmas}
\makeatletter
\newlabel{lemma-iso-global-hom}{{0B6I}{source}{}{sourcepointer-7}{}}
\newlabel{lemma-iso-global-hom@cref}{{[originallemma][0][]0B6I}{[1][0][]source}{}{}{}}
\makeatother
\crefname{originallemma}{original lemma}{original lemmas}
\Crefname{originallemma}{Original lemma}{Original lemmas}
\makeatletter
\newlabel{perfect-lemma-compare-base-change}{{08IB}{source}{}{sourcepointer-8}{}}
\newlabel{perfect-lemma-compare-base-change@cref}{{[originallemma][0][]08IB}{[1][0][]source}{}{}{}}
\makeatother
\crefname{originallemma}{original lemma}{original lemmas}
\Crefname{originallemma}{Original lemma}{Original lemmas}
\makeatletter
\newlabel{perfect-lemma-quasi-coherence-direct-image}{{08D5}{source}{}{sourcepointer-9}{}}
\newlabel{perfect-lemma-quasi-coherence-direct-image@cref}{{[originallemma][0][]08D5}{[1][0][]source}{}{}{}}
\makeatother
\crefname{originalequation}{original equation}{original equations}
\Crefname{originalequation}{Original equation}{Original equations}
\makeatletter
\newlabel{more-algebra-equation-base-change-RHom}{{0E1X}{source}{}{sourcepointer-10}{}}
\newlabel{more-algebra-equation-base-change-RHom@cref}{{[originalequation][0][]0E1X}{[1][0][]source}{}{}{}}
\makeatother
\crefname{originallemma}{original lemma}{original lemmas}
\Crefname{originallemma}{Original lemma}{Original lemmas}
\makeatletter
\newlabel{perfect-lemma-affine-compare-bounded}{{06Z0}{source}{}{sourcepointer-11}{}}
\newlabel{perfect-lemma-affine-compare-bounded@cref}{{[originallemma][0][]06Z0}{[1][0][]source}{}{}{}}
\makeatother
\crefname{originallemma}{original lemma}{original lemmas}
\Crefname{originallemma}{Original lemma}{Original lemmas}
\makeatletter
\newlabel{perfect-lemma-quasi-coherence-internal-hom}{{0A6H}{source}{}{sourcepointer-12}{}}
\newlabel{perfect-lemma-quasi-coherence-internal-hom@cref}{{[originallemma][0][]0A6H}{[1][0][]source}{}{}{}}
\makeatother
\crefname{originallemma}{original lemma}{original lemmas}
\Crefname{originallemma}{Original lemma}{Original lemmas}
\makeatletter
\newlabel{perfect-lemma-flat-proper-perfect-direct-image-general}{{0B91}{source}{}{sourcepointer-13}{}}
\newlabel{perfect-lemma-flat-proper-perfect-direct-image-general@cref}{{[originallemma][0][]0B91}{[1][0][]source}{}{}{}}
\makeatother
\crefname{originallemma}{original lemma}{original lemmas}
\Crefname{originallemma}{Original lemma}{Original lemmas}
\makeatletter
\newlabel{more-algebra-lemma-base-change-RHom}{{0A6A}{source}{}{sourcepointer-14}{}}
\newlabel{more-algebra-lemma-base-change-RHom@cref}{{[originallemma][0][]0A6A}{[1][0][]source}{}{}{}}
\makeatother
\crefname{originallemma}{original lemma}{original lemmas}
\Crefname{originallemma}{Original lemma}{Original lemmas}
\makeatletter
\newlabel{more-morphisms-lemma-perfect-proper-perfect-direct-image}{{0B6G}{source}{}{sourcepointer-15}{}}
\newlabel{more-morphisms-lemma-perfect-proper-perfect-direct-image@cref}{{[originallemma][0][]0B6G}{[1][0][]source}{}{}{}}
\makeatother
\crefname{originallemma}{original lemma}{original lemmas}
\Crefname{originallemma}{Original lemma}{Original lemmas}
\makeatletter
\newlabel{perfect-lemma-direct-image-coherent}{{08E2}{source}{}{sourcepointer-16}{}}
\newlabel{perfect-lemma-direct-image-coherent@cref}{{[originallemma][0][]08E2}{[1][0][]source}{}{}{}}
\makeatother
\crefname{originallemma}{original lemma}{original lemmas}
\Crefname{originallemma}{Original lemma}{Original lemmas}
\makeatletter
\newlabel{perfect-lemma-identify-pseudo-coherent-noetherian}{{08E8}{source}{}{sourcepointer-17}{}}
\newlabel{perfect-lemma-identify-pseudo-coherent-noetherian@cref}{{[originallemma][0][]08E8}{[1][0][]source}{}{}{}}
\makeatother
\crefname{originallemma}{original lemma}{original lemmas}
\Crefname{originallemma}{Original lemma}{Original lemmas}
\makeatletter
\newlabel{lemma-flat-precompose-pus}{{0A9K}{source}{}{sourcepointer-18}{}}
\newlabel{lemma-flat-precompose-pus@cref}{{[originallemma][0][]0A9K}{[1][0][]source}{}{}{}}
\makeatother
\crefname{originalreference}{original reference}{original references}
\Crefname{originalreference}{Original reference}{Original references}
\makeatletter
\newlabel{perfect-remark-multiplication-map}{{0BZA}{source}{}{sourcepointer-19}{}}
\newlabel{perfect-remark-multiplication-map@cref}{{[originalreference][0][]0BZA}{[1][0][]source}{}{}{}}
\makeatother
\crefname{originalsection}{original section}{original sections}
\Crefname{originalsection}{Original section}{Original sections}
\makeatletter
\newlabel{more-algebra-section-base-change-RHom}{{0E1V}{source}{}{sourcepointer-20}{}}
\newlabel{more-algebra-section-base-change-RHom@cref}{{[originalsection][0][]0E1V}{[1][0][]source}{}{}{}}
\makeatother
\title{Base change for the right adjoint of pushforward under Tor independence, in three cases}
\author{The Stacks Project Authors}
\date{stacks-project:duality.tex:local-snapshot-20261012; GNU FDL 1.2 or later}
\begin{document}
\maketitle
\noindent\textit{Editorial note.} All delivered mathematical text is the author's own, copied from the pinned Stacks Project chapter file duality.tex; no statement or proof step was rewritten, reordered or removed.\par\medskip
\noindent\textit{Editorial note.} Preamble provenance: the pinned chapter file opens with an input of the shared preamble and carries no preamble text of its own, so the project's single Stacks-wide preamble staged at sources/stacks-duality/preamble.tex is excerpted verbatim under the recipe fixed by LANE\_BRIEF 22.1, as three spans, lines 8-13 (verbatim package and the reference/slogan/history environments), lines 14-18 (Xy-pic setup) and lines 154-217 (the native theorem declarations and mathematical macros). Nothing was added, removed or redefined; the wrapper class and the amsmath/amssymb/amsthm/graphicx base packages are supplied by the assembler.\par\medskip
\noindent\textit{Editorial note.} The counted statement and proof contain no cross reference, no citation and no included file, so the delivered body is self contained.\par\medskip
\noindent\textit{Editorial note.} Canonical location of the counted proposition: Stacks tag 0AA8 (sources/stacks-duality/pinned-original-source.tex line 1636).\par\medskip
\noindent\textit{Allowed standard prerequisites (editorial).}\begin{itemize}
\item The right adjoint a of pushforward, the base change diagram and its base change map. Same chapter, diagram labelled equation-base-change and equation labelled equation-base-change-map, invoked in the statement, and Lemma labelled lemma-iso-global-hom, cited at proof lines 1693-1697 and 1699-1703. Fixes the two functors and identifies both global Hom complexes with Hom on the base.
\item Perfect generators of the derived category of quasi-coherent sheaves. Derived Categories of Schemes, Theorem labelled perfect-theorem-bondal-van-den-Bergh and Remark labelled perfect-remark-pullback-generator, cited at proof lines 1661-1666. Reduces the isomorphism of objects to the isomorphism of Hom complexes out of a generator.
\item Base change for Hom out of a perfect object and comparison of the pushforwards. Derived Categories of Schemes, Lemmas labelled perfect-lemma-affine-morphism-and-hom-out-of-perfect and perfect-lemma-compare-base-change, cited at proof lines 1687-1691 and 1703-1707. Computes the left hand side and rewrites the right hand side.
\item Base change for Hom over a ring, in the four cases. More on Algebra, Equation labelled more-algebra-equation-base-change-RHom, Section labelled more-algebra-section-base-change-RHom and Lemma labelled more-algebra-lemma-base-change-RHom parts (1) and (4), cited at proof lines 1713-1731 and 1819-1823. Says when beta is an isomorphism and identifies it with the map whose commutativity is proved at the end.
\item The invertible transformation between the two compositions of functors. Same chapter, Lemma labelled lemma-flat-precompose-pus, cited at proof lines 1769-1773, and Derived Categories of Schemes, Remark labelled perfect-remark-multiplication-map, cited at proof lines 1787-1791. Supplies the maps mu and mu' whose commutativity reduces the final check to a triangle.
\end{itemize}\medskip
\ifdefined\hypertarget\hypertarget{sourcepointer-1}{}\fi
\noindent\textit{Source reference.} Original equation 0A9J: sources/stacks-duality/pinned-original-source.tex, source lines 927--927. Same-chapter pointer into the author environment equation at source line 927; invoked at proof line None of the counted proof. It is retained only as a navigation pointer; every load-bearing step is inside the delivered span.\par\medskip
\ifdefined\hypertarget\hypertarget{sourcepointer-2}{}\fi
\noindent\textit{Source reference.} Original equation 0AA6: sources/stacks-duality/pinned-original-source.tex, source lines 1192--1192. Same-chapter pointer into the author environment equation at source line 1192; invoked at proof line 1666 of the counted proof. It is retained only as a navigation pointer; every load-bearing step is inside the delivered span.\par\medskip
\ifdefined\hypertarget\hypertarget{sourcepointer-3}{}\fi
\noindent\textit{Source reference.} Original theorem 09IS: sources/stacks-chow/tags.txt, source lines 12050--12050. Cross-chapter standard prerequisite located in the project tag list staged in the local snapshot; invoked at proof line 1663 of the counted proof. It is retained only as a navigation pointer; every load-bearing step is inside the delivered span.\par\medskip
\ifdefined\hypertarget\hypertarget{sourcepointer-4}{}\fi
\noindent\textit{Source reference.} Original reference 0BQT: sources/stacks-chow/tags.txt, source lines 14746--14746. Cross-chapter standard prerequisite located in the project tag list staged in the local snapshot; invoked at proof line 1665 of the counted proof. It is retained only as a navigation pointer; every load-bearing step is inside the delivered span.\par\medskip
\ifdefined\hypertarget\hypertarget{sourcepointer-5}{}\fi
\noindent\textit{Source reference.} Original section 0B6A: sources/stacks-chow/tags.txt, source lines 14063--14063. Cross-chapter standard prerequisite located in the project tag list staged in the local snapshot; invoked at proof line 1675 of the counted proof. It is retained only as a navigation pointer; every load-bearing step is inside the delivered span.\par\medskip
\ifdefined\hypertarget\hypertarget{sourcepointer-6}{}\fi
\noindent\textit{Source reference.} Original lemma 0AA7: sources/stacks-chow/tags.txt, source lines 12975--12975. Cross-chapter standard prerequisite located in the project tag list staged in the local snapshot; invoked at proof line 1690 of the counted proof. It is retained only as a navigation pointer; every load-bearing step is inside the delivered span.\par\medskip
\ifdefined\hypertarget\hypertarget{sourcepointer-7}{}\fi
\noindent\textit{Source reference.} Original lemma 0B6I: sources/stacks-duality/pinned-original-source.tex, source lines 878--878. Same-chapter pointer into the author environment lemma at source line 878; invoked at proof line 1693 of the counted proof. It is retained only as a navigation pointer; every load-bearing step is inside the delivered span.\par\medskip
\ifdefined\hypertarget\hypertarget{sourcepointer-8}{}\fi
\noindent\textit{Source reference.} Original lemma 08IB: sources/stacks-chow/tags.txt, source lines 10809--10809. Cross-chapter standard prerequisite located in the project tag list staged in the local snapshot; invoked at proof line 1704 of the counted proof. It is retained only as a navigation pointer; every load-bearing step is inside the delivered span.\par\medskip
\ifdefined\hypertarget\hypertarget{sourcepointer-9}{}\fi
\noindent\textit{Source reference.} Original lemma 08D5: sources/stacks-chow/tags.txt, source lines 10625--10625. Cross-chapter standard prerequisite located in the project tag list staged in the local snapshot; invoked at proof line 1708 of the counted proof. It is retained only as a navigation pointer; every load-bearing step is inside the delivered span.\par\medskip
\ifdefined\hypertarget\hypertarget{sourcepointer-10}{}\fi
\noindent\textit{Source reference.} Original equation 0E1X: sources/stacks-chow/tags.txt, source lines 17587--17587. Cross-chapter standard prerequisite located in the project tag list staged in the local snapshot; invoked at proof line 1717 of the counted proof. It is retained only as a navigation pointer; every load-bearing step is inside the delivered span.\par\medskip
\ifdefined\hypertarget\hypertarget{sourcepointer-11}{}\fi
\noindent\textit{Source reference.} Original lemma 06Z0: sources/stacks-chow/tags.txt, source lines 8904--8904. Cross-chapter standard prerequisite located in the project tag list staged in the local snapshot; invoked at proof line 1719 of the counted proof. It is retained only as a navigation pointer; every load-bearing step is inside the delivered span.\par\medskip
\ifdefined\hypertarget\hypertarget{sourcepointer-12}{}\fi
\noindent\textit{Source reference.} Original lemma 0A6H: sources/stacks-chow/tags.txt, source lines 12845--12845. Cross-chapter standard prerequisite located in the project tag list staged in the local snapshot; invoked at proof line 1720 of the counted proof. It is retained only as a navigation pointer; every load-bearing step is inside the delivered span.\par\medskip
\ifdefined\hypertarget\hypertarget{sourcepointer-13}{}\fi
\noindent\textit{Source reference.} Original lemma 0B91: sources/stacks-chow/tags.txt, source lines 14159--14159. Cross-chapter standard prerequisite located in the project tag list staged in the local snapshot; invoked at proof line 1725 of the counted proof. It is retained only as a navigation pointer; every load-bearing step is inside the delivered span.\par\medskip
\ifdefined\hypertarget\hypertarget{sourcepointer-14}{}\fi
\noindent\textit{Source reference.} Original lemma 0A6A: sources/stacks-chow/tags.txt, source lines 12838--12838. Cross-chapter standard prerequisite located in the project tag list staged in the local snapshot; invoked at proof line 1727 of the counted proof. It is retained only as a navigation pointer; every load-bearing step is inside the delivered span.\par\medskip
\ifdefined\hypertarget\hypertarget{sourcepointer-15}{}\fi
\noindent\textit{Source reference.} Original lemma 0B6G: sources/stacks-chow/tags.txt, source lines 14069--14069. Cross-chapter standard prerequisite located in the project tag list staged in the local snapshot; invoked at proof line 1730 of the counted proof. It is retained only as a navigation pointer; every load-bearing step is inside the delivered span.\par\medskip
\ifdefined\hypertarget\hypertarget{sourcepointer-16}{}\fi
\noindent\textit{Source reference.} Original lemma 08E2: sources/stacks-chow/tags.txt, source lines 10658--10658. Cross-chapter standard prerequisite located in the project tag list staged in the local snapshot; invoked at proof line 1735 of the counted proof. It is retained only as a navigation pointer; every load-bearing step is inside the delivered span.\par\medskip
\ifdefined\hypertarget\hypertarget{sourcepointer-17}{}\fi
\noindent\textit{Source reference.} Original lemma 08E8: sources/stacks-chow/tags.txt, source lines 10664--10664. Cross-chapter standard prerequisite located in the project tag list staged in the local snapshot; invoked at proof line 1736 of the counted proof. It is retained only as a navigation pointer; every load-bearing step is inside the delivered span.\par\medskip
\ifdefined\hypertarget\hypertarget{sourcepointer-18}{}\fi
\noindent\textit{Source reference.} Original lemma 0A9K: sources/stacks-duality/pinned-original-source.tex, source lines 957--957. Same-chapter pointer into the author environment lemma at source line 957; invoked at proof line 1774 of the counted proof. It is retained only as a navigation pointer; every load-bearing step is inside the delivered span.\par\medskip
\ifdefined\hypertarget\hypertarget{sourcepointer-19}{}\fi
\noindent\textit{Source reference.} Original reference 0BZA: sources/stacks-chow/tags.txt, source lines 15043--15043. Cross-chapter standard prerequisite located in the project tag list staged in the local snapshot; invoked at proof line 1789 of the counted proof. It is retained only as a navigation pointer; every load-bearing step is inside the delivered span.\par\medskip
\ifdefined\hypertarget\hypertarget{sourcepointer-20}{}\fi
\noindent\textit{Source reference.} Original section 0E1V: sources/stacks-chow/tags.txt, source lines 17585--17585. Cross-chapter standard prerequisite located in the project tag list staged in the local snapshot; invoked at proof line 1821 of the counted proof. It is retained only as a navigation pointer; every load-bearing step is inside the delivered span.\par\medskip
\begin{lemma}
\label{lemma-more-base-change}
In diagram (\ref{equation-base-change}) assume
\begin{enumerate}
\item $g : Y' \to Y$ is a morphism of affine schemes,
\item $f : X \to Y$ is proper, and
\item $f$ and $g$ are Tor independent.
\end{enumerate}
Then the base change map (\ref{equation-base-change-map}) induces an
isomorphism
$$
L(g')^*a(K) \longrightarrow a'(Lg^*K)
$$
in the following cases
\begin{enumerate}
\item for all $K \in D_\QCoh(\mathcal{O}_Y)$ if $f$
is flat of finite presentation,
\item for all $K \in D_\QCoh(\mathcal{O}_Y)$ if $f$
is perfect and $Y$ Noetherian,
\item for $K \in D_\QCoh^+(\mathcal{O}_Y)$ if $g$ has finite Tor dimension
and $Y$ Noetherian.
\end{enumerate}
\end{lemma}

\begin{proof}
Write $Y = \Spec(A)$ and $Y' = \Spec(A')$. As a base change of an affine
morphism, the morphism $g'$ is affine. Let $M$ be a perfect generator
for $D_\QCoh(\mathcal{O}_X)$, see Derived Categories of Schemes, Theorem
\ref{perfect-theorem-bondal-van-den-Bergh}. Then $L(g')^*M$ is a
generator for $D_\QCoh(\mathcal{O}_{X'})$, see
Derived Categories of Schemes, Remark \ref{perfect-remark-pullback-generator}.
Hence it suffices to show that (\ref{equation-base-change-map})
induces an isomorphism
\begin{equation}
\label{equation-iso}
R\Hom_{X'}(L(g')^*M, L(g')^*a(K))
\longrightarrow
R\Hom_{X'}(L(g')^*M, a'(Lg^*K))
\end{equation}
of global hom complexes, see
Cohomology, Section \ref{cohomology-section-global-RHom},
as this will imply the cone of $L(g')^*a(K) \to a'(Lg^*K)$
is zero.
The structure of the proof is as follows: we will first show that
these Hom complexes are isomorphic and in the last part of the proof
we will show that the isomorphism is induced by (\ref{equation-iso}).

\medskip\noindent
The left hand side. Because $M$ is perfect, the canonical map
$$
R\Hom_X(M, a(K)) \otimes^\mathbf{L}_A A'
\longrightarrow
R\Hom_{X'}(L(g')^*M, L(g')^*a(K))
$$
is an isomorphism by Derived Categories of Schemes, Lemma
\ref{perfect-lemma-affine-morphism-and-hom-out-of-perfect}.
We can combine this with the isomorphism
$R\Hom_Y(Rf_*M, K) = R\Hom_X(M, a(K))$
of Lemma \ref{lemma-iso-global-hom}
to get that the left hand side equals
$R\Hom_Y(Rf_*M, K) \otimes^\mathbf{L}_A A'$.

\medskip\noindent
The right hand side. Here we first use the isomorphism
$$
R\Hom_{X'}(L(g')^*M, a'(Lg^*K)) = R\Hom_{Y'}(Rf'_*L(g')^*M, Lg^*K)
$$
of Lemma \ref{lemma-iso-global-hom}. Then we use the base change
map $Lg^*Rf_*M \to Rf'_*L(g')^*M$ is an isomorphism by
Derived Categories of Schemes, Lemma \ref{perfect-lemma-compare-base-change}.
Hence we may rewrite this as $R\Hom_{Y'}(Lg^*Rf_*M, Lg^*K)$.
Since $Y$, $Y'$ are affine and $K$, $Rf_*M$ are in $D_\QCoh(\mathcal{O}_Y)$
(Derived Categories of Schemes, Lemma
\ref{perfect-lemma-quasi-coherence-direct-image})
we have a canonical map
$$
\beta :
R\Hom_Y(Rf_*M, K) \otimes^\mathbf{L}_A A'
\longrightarrow
R\Hom_{Y'}(Lg^*Rf_*M, Lg^*K)
$$
in $D(A')$. This is the arrow
More on Algebra, Equation (\ref{more-algebra-equation-base-change-RHom})
where we have used Derived Categories of Schemes, Lemmas
\ref{perfect-lemma-affine-compare-bounded} and
\ref{perfect-lemma-quasi-coherence-internal-hom}
to translate back and forth into algebra.
\begin{enumerate}
\item If $f$ is flat and of finite presentation, the complex $Rf_*M$
is perfect on $Y$ by Derived Categories of Schemes, Lemma
\ref{perfect-lemma-flat-proper-perfect-direct-image-general}
and $\beta$ is an isomorphism by
More on Algebra, Lemma \ref{more-algebra-lemma-base-change-RHom} part (1).
\item If $f$ is perfect and $Y$ Noetherian, the complex $Rf_*M$
is perfect on $Y$ by More on Morphisms, Lemma
\ref{more-morphisms-lemma-perfect-proper-perfect-direct-image}
and $\beta$ is an isomorphism as before.
\item If $g$ has finite tor dimension and $Y$ is Noetherian,
the complex $Rf_*M$ is pseudo-coherent on $Y$
(Derived Categories of Schemes, Lemmas
\ref{perfect-lemma-direct-image-coherent} and
\ref{perfect-lemma-identify-pseudo-coherent-noetherian})
and $\beta$ is an isomorphism by
More on Algebra, Lemma \ref{more-algebra-lemma-base-change-RHom} part (4).
\end{enumerate}
We conclude that we obtain the same answer as in the previous paragraph.

\medskip\noindent
In the rest of the proof we show that the identifications of
the left and right hand side of (\ref{equation-iso})
given in the second and third paragraph are in fact given by
(\ref{equation-iso}). To make our formulas manageable
we will use $(-, -)_X = R\Hom_X(-, -)$, use $- \otimes A'$
in stead of $- \otimes_A^\mathbf{L} A'$, and we will abbreviate
$g^* = Lg^*$ and $f_* = Rf_*$. Consider the following
commutative diagram
$$
\xymatrix{
((g')^*M, (g')^*a(K))_{X'} \ar[d] &
(M, a(K))_X \otimes A' \ar[l]^-\alpha \ar[d] &
(f_*M, K)_Y \otimes A' \ar@{=}[l] \ar[d] \\
((g')^*M, (g')^*a(g_*g^*K))_{X'} &
(M, a(g_*g^*K))_X \otimes A' \ar[l]^-\alpha &
(f_*M, g_*g^*K)_Y \otimes A' \ar@{=}[l] \ar@/_4pc/[dd]_{\mu'} \\
((g')^*M, (g')^*g'_*a'(g^*K))_{X'} \ar[u] \ar[d] &
(M, g'_*a'(g^*K))_X \otimes A' \ar[u] \ar[l]^-\alpha \ar[ld]^\mu &
(f_*M, K) \otimes A' \ar[d]^\beta \\
((g')^*M, a'(g^*K))_{X'} &
(f'_*(g')^*M, g^*K)_{Y'} \ar@{=}[l] \ar[r] &
(g^*f_*M, g^*K)_{Y'}
}
$$
The arrows labeled $\alpha$ are the maps from
Derived Categories of Schemes, Lemma
\ref{perfect-lemma-affine-morphism-and-hom-out-of-perfect}
for the diagram with corners $X', X, Y', Y$.
The upper part of the diagram is commutative as the horizontal arrows are
functorial in the entries.
The middle vertical arrows come from the invertible transformation
$g'_* \circ a' \to a \circ g_*$  of Lemma \ref{lemma-flat-precompose-pus}
and therefore the middle square is commutative.
Going down the left hand side is (\ref{equation-iso}).
The upper horizontal arrows provide the identifications used in the
second paragraph of the proof.
The lower horizontal arrows including $\beta$ provide the identifications
used in the third paragraph of the proof. Given $E \in D(A)$,
$E' \in D(A')$, and $c : E \to E'$ in $D(A)$ we will denote
$\mu_c : E \otimes A' \to E'$ the map induced by $c$
and the adjointness of restriction and base change;
if $c$ is clear we write $\mu = \mu_c$, i.e., we
drop $c$ from the notation. The map $\mu$ in the diagram is of this
form with $c$ given by the identification
$(M, g'_*a(g^*K))_X = ((g')^*M, a'(g^*K))_{X'}$
; the triangle involving $\mu$ is commutative by
Derived Categories of Schemes, Remark \ref{perfect-remark-multiplication-map}.

\medskip\noindent
Observe that
$$
\xymatrix{
(M, a(g_*g^*K))_X &
(f_*M, g_* g^*K)_Y \ar@{=}[l] &
(g^*f_*M, g^*K)_{Y'} \ar@{=}[l] \\
(M, g'_* a'(g^*K))_X \ar[u] &
((g')^*M, a'(g^*K))_{X'} \ar@{=}[l] &
(f'_*(g')^*M, g^*K)_{Y'} \ar@{=}[l] \ar[u]
}
$$
is commutative by the very definition of the transformation
$g'_* \circ a' \to a \circ g_*$. Letting $\mu'$ be as above
corresponding to the identification
$(f_*M, g_*g^*K)_X = (g^*f_*M, g^*K)_{Y'}$, then the
hexagon commutes as well. Thus it suffices to show that
$\beta$ is equal to the composition of
$(f_*M, K)_Y \otimes A' \to (f_*M, g_*g^*K)_X \otimes A'$
and $\mu'$. To do this, it suffices to prove the two induced maps
$(f_*M, K)_Y \to (g^*f_*M, g^*K)_{Y'}$ are the same.
In other words, it suffices to show the diagram
$$
\xymatrix{
R\Hom_A(E, K) \ar[rr]_{\text{induced by }\beta} \ar[rd] & &
R\Hom_{A'}(E \otimes_A^\mathbf{L} A', K \otimes_A^\mathbf{L} A') \\
& R\Hom_A(E, K \otimes_A^\mathbf{L} A') \ar[ru]
}
$$
commutes for all $E, K \in D(A)$. Since this is how $\beta$ is constructed in
More on Algebra, Section \ref{more-algebra-section-base-change-RHom}
the proof is complete.
\end{proof}

\end{document}
